\section{Quantitative consequences and comparison of regimes}
\label{sec:consequences}

This appendix records the dependence of the estimates on an external
exponential normalization and compares the static and parabolic scales.

\begin{corollary}[A sufficient coefficient at the growing-mode boundary]
\label{cor:critical-coefficient}
Keep the cutoffs and constants of Theorem~\ref{thm:ub-flux} fixed.
Fix \(p_0,\gamma>0\), choose \(R(L)>0\) so that its geometric
hypotheses hold for all sufficiently large \(L\), let \(a=p_0L\) and
\(\omega=\exp(-\gamma L^2)\), and use the normalized quantities
\[
 P=R^{-2}\omega M^{\rm plat}_{\boldsymbol n,R},
 \qquad X=(\omega/R)[\mathcal T_{\boldsymbol n,R}]_+.
\]
If \(q(L)\ge1\) is integer-valued and
\(q(L)/L^{5/2}\to\kappa\in[0,\infty)\), the coefficient
\[
 \mathcal B_L
 =C a^{2/3}q(L)^7\omega^{1/3}
          \exp(12\sqrt2\,q(L)\sqrt{\Delta_a})
\]
in \(X\le\mathcal B_LP^{2/3}\) satisfies
\begin{equation}\label{eq:critical-rate}
 \lim_{L\to\infty}\frac{\log\mathcal B_L}{L^2}
 =-\frac{\gamma}{3}
       +12\sqrt2\,\kappa\sqrt{\frac{\delta+\delta_0}{p_0}}.
\end{equation}
In particular, the coefficient still decays at a strictly negative
quadratic logarithmic rate whenever
\begin{equation}\label{eq:critical-kappa}
 \kappa<
 \frac{\gamma}{36\sqrt2}
           \sqrt{\frac{p_0}{\delta+\delta_0}}.
\end{equation}
The same conclusion holds with a strict limsup in
\eqref{eq:critical-kappa}.
\end{corollary}

\begin{proof}
Theorem~\ref{thm:ub-flux} gives the normalized inequality by exact
cancellation of the powers of \(R\):
\[
 \frac{\omega}{R}\,
     a^{2/3}R^{-1/3}(R^2\omega^{-1}P)^{2/3}
 =a^{2/3}\omega^{1/3}P^{2/3}.
\]
Since \(\Delta_a=(\delta+\delta_0)/(p_0L-\delta_0)\),
\[
 \frac{12\sqrt2\,q(L)\sqrt{\Delta_a}}{L^2}
 =12\sqrt2\,\frac{q(L)}{L^{5/2}}
       \sqrt{\frac{\delta+\delta_0}{p_0-\delta_0/L}}.
\]
The limit of this expression is the positive term in
\eqref{eq:critical-rate}. The logarithms of \(C,a^{2/3}\), and \(q^7\),
divided by \(L^2\), tend to zero, and the weight contributes
\(-\gamma/3\). The same calculation with a limsup proves the last
assertion.
\end{proof}

This is a sufficient range for the proved upper-bound coefficient.
Neither equality in \eqref{eq:critical-kappa} nor its failure establishes
a change in the actual flux. The numerical constants come from the cited
extrapolation input. The spatial witnesses in
Section~\ref{sec:sharpness} do not supply a matching time-integrated
lower bound for this growing-mode boundary.

For the scale choice
\(p_0=32/63\), \(\gamma=2/3969\), and
\(R=\exp(-9L^2/40000)\), the subcritical rate in
Remark~\ref{rem:ub-window} is \(-2/11907\). The exact Chebyshev family
of Theorem~\ref{thm:parabolic-shear} has the same normalized rate in
its smaller regime \(m=o(A^2)\). These two statements have different
strengths: the first is a uniform upper-bound coefficient; the second
is an attained rate for the explicitly chosen family.

A power-law parametrization makes the distinction between the static
and parabolic scales visible. Let \(m\) be the nearest positive even
integer to \(A^\sigma\), with fixed \(1/2<\sigma<2\). Then
\begin{equation}\label{eq:power-law-regimes}
 \frac{m}{\sqrt A}\sim A^{\sigma-1/2},\qquad
 \mu\sim A^{(2\sigma-1)/3},\qquad
 \mu^2\sim A^{(4\sigma-2)/3}.
\end{equation}
The first scale describes the logarithm of a static edge gain.
The second describes the edge-dependent exponential factor in the
complete-clock quotient, with its algebraic prefactors retained.
The third describes the common heat amplitude that cancels from the
normalized comparison. In this regime \(\mu=o(A)\), and in particular
\(\mu=o(L^2)\) when \(A\) is comparable with \(L\).

The polynomial factors in the flux quotient should not be discarded.
For example, \(m/\sqrt A\to\infty\) allows \(\mu\) to grow more slowly
than \(\log A\). In that case an expression described only as
\(\exp(\Theta(\mu))\) need not imply that the unweighted quotient
diverges. The two-sided inequality in
Theorem~\ref{thm:parabolic-shear}, rather than that shorthand, is the
statement used here.
